% CLS: incoming finite heat/translation applied to actual CC coefficient sheaves.
\section{Heat and translation projectors on the localized adelic coefficient family}
\label{sec:cc-localized-shift}

\paragraph{Provenance and original objects.}
The incoming programme source is \emph{Dynamical Synchronization, Heat
Remapping, and the Full Supported Cauchy--Weil Form}, 23 September 2026,
\nolinkurl{RECONSTRUCTION.tex}, ST1--12 and SY1--16. Its archive is retained
unchanged with SHA-256
\nolinkurl{70810361fd8fd66f9cc7154ed919bc5c0c15d003e3ec72504b173bbb24aeed15}.
Those finite-source sections were read in full. Their identities are
proved below in the exact receiving sheaves; neither their later
analytic claims nor the sign of a Weil form is assumed. The adelic
coefficient input is Alain Connes and Caterina Consani's
\href{https://arxiv.org/abs/2501.06560v1}{\emph{Knots, primes and class
field theory}}, Section 4, extended with full proofs in CDV1--25 and
ULC1--18. The original arithmetic function remains $\zeta$, with
the exact factors of ULC18 and ZTC10. Finite formal-coordinate heat
here is not an assertion that those coefficients follow Newman heat.

\subsection{Finite matrices on the original semimodule sheaf}
Use the actual localized sheaf $\mathcal M_k$ of ULC4 on
$Y=\operatorname{Spec}G_L(\mathbb Z)$ and fix $N\ge2$.
Its independent coefficient source is $\mathcal M_k^{N+1}$,
ordered by $1,Z,\ldots,Z^N$. A constant split complex scalar
$(c,\lambda)\in G_L(\mathbb C)$ acts on arithmetic-open sections by
$(w,\mu)\mapsto(cw,\lambda\wedge\mu)$ and on a support stalk
$L_J$ by meeting with $[\lambda]_J$. These maps commute with
restrictions by ULC4 and with the $\mathcal O_Y$ action. Nonzero
rational coefficients, including their factorial denominators,
therefore define the stated endomorphisms without treating them
as regular rational functions at every arithmetic prime.

Let $D_{i,i+1}=(i+1)^\bullet$ and every structural zero be $\tau$.
For split parameters $a,b$ define
\begin{equation}\tag{CLS1}
 F_{a,b}=\sum_{n=0}^{N}\frac{(-aD^2+bD)^n}{n!},\qquad
 (F_{a,b})_{ij}=\sum_{q=0}^{\lfloor (j-i)/2\rfloor}
 \left(\frac{(-1)^qj!}{i!q!(j-i-2q)!}\right)^\bullet
 a^qb^{j-i-2q}\quad(j\ge i).
\end{equation}
Below the diagonal the entries are $\tau$. In the first formula
the reciprocal factorial is fully supported, and the zeroth
power is the free identity matrix. Each generator strictly lowers
degree, so all omitted higher terms are the global zero matrix.
To obtain the second formula, select $q$ factors of $-aD^2$ and
$j-i-2q$ factors of $bD$, then use the binomial coefficient and
$D^{j-i}(Z^j)=j!Z^i/i!$. This keeps every sign and denominator.
Commuting the generators and using the finite binomial identity gives
\begin{equation}\tag{CLS2}
 F_{a,b}F_{c,d}=F_{a+c,b+d},\qquad F_{\tau,\tau}=I.
\end{equation}
This identity is valid in the semiring of matrices: the binomial
proof uses distributivity, not subtraction of supported zeros.

At $a=(0,\lambda)$, $b=(0,\mu)$, every positive-degree amplitude
vanishes but its original label remains. An odd degree drop has
label $\mu$ because every term contains $b$ and the term $q=0$
has exactly that label. A positive even drop has the endpoint
labels $\lambda$ and $\mu$, with all intermediate labels their
meet, so its resulting label is $\lambda\vee\mu$. Thus
\begin{equation}\tag{CLS3}
 P=F_{e,\tau},\quad Q=F_{e,e}=F_{\tau,e},\quad
 P^2=P,\qquad Q^2=Q,\qquad PQ=QP=Q.
\end{equation}
At an arithmetic stalk, these maps keep each amplitude and change
the independent labels by
\begin{equation}\tag{CLS4}
 (P\boldsymbol\lambda)_i=
       \bigvee_{j\ge i,\ j-i\ \mathrm{even}}\lambda_j,
 \qquad (Q\boldsymbol\lambda)_i=\bigvee_{j\ge i}\lambda_j.
\end{equation}
This follows directly from matrix multiplication: the off-diagonal
entries are zero-amplitude elements, so only the diagonal term
contributes an amplitude, while joins record all admissible labels.
Their images have respectively nested same-parity labels and a
single nested flag. They are sheaf retractions onto their images.

\subsection{Localization preserves the projectors' difference}
At a support prime $J$, ULC8 sends every nonzero rational scalar
and $e$ to $1\in L_J$, while $\tau$ remains zero. Consequently
the localized $P_J$ and $Q_J$ are exactly the join operators CLS4
on $(L_J)^{N+1}$, now with their localized labels. They are
different at every support prime: the input with entry $1$ only
at degree one and zero elsewhere has output degree zero equal
to $0$ under $P_J$ and $1$ under $Q_J$. A proper prime localization
has $0\ne1$: an equality would require an element outside $J$
to meet $0$ and $1$ equally, forcing that element to be $0\in J$.
Thus the example is nonzero in every support stalk.

The additive group reflection instead sends both projectors to
the identity on $(i_*\mathcal H_k)^{N+1}$: every off-diagonal
entry in them has amplitude zero. It follows, by applying the
explicit CDV complex coordinatewise, that
\begin{equation}\tag{CLS5}
 H^j(P)=H^j(Q)=\mathrm{id}\quad\text{on }
 H^j(X,\mathcal H_k)^{N+1},\qquad j\ge0,
\end{equation}
where these cohomology maps are those of the amplitude reflection.
Every mixed-prime summand of CDV13 is included. The equality is
not a claim that $P=Q$ on $\mathcal M_k^{N+1}$; their explicit
support-stalk discrepancy proves the contrary. This gives the
actual location of the information lost by ordinary additive
cohomology in this calculation.

For supported $t,b\in\mathbb C$, CLS2 gives
\begin{equation}\tag{CLS6}
 F_{t^\bullet,b^\bullet}F_{(-t)^\bullet,(-b)^\bullet}=Q.
\end{equation}
Also $FQ=QF=F$, since adding $e$ to a supported parameter leaves
that parameter unchanged. Hence these maps form a group on
$Q\mathcal M_k^{N+1}$ with identity $Q$. On a support stalk all
these supported-parameter matrices have the same image $Q_J$:
in CLS1 every positive degree has at least one top-labelled term,
and localization retains its label even when its amplitude cancels.
They therefore act as the identity on the localized projective image.
On the full free localized source, their common action is the
nonidentity projector $Q_J$. The arithmetic amplitude group is the
usual invertible polynomial heat/translation action.

\subsection{The common-label comparison and its exact fibres}
Define $\mathcal M_k[Z]_{\le N}^{\mathrm{one}}$ by the construction
ULC4 applied to the vector sheaf $\mathcal H_k[Z]_{\le N}$.
Its arithmetic-open sections have one common label for the entire
polynomial; its support restriction is $\mathcal O_S$.
There is a surjective sheaf map
\begin{equation}\tag{CLS7}
 \eta_N:\mathcal M_k^{N+1}\longrightarrow
 \mathcal M_k[Z]_{\le N}^{\mathrm{one}},\qquad
 ((w_i,\lambda_i))\longmapsto
       \left(\sum_iw_iZ^i,\bigvee_i\lambda_i\right).
\end{equation}
On support opens it is the join of the coordinate sections.
A nonzero polynomial forces at least one nonzero coefficient,
hence top label; so the image belongs to the original carrier.
Distributivity proves compatibility with scalar meets, and joins
prove additivity. A polynomial with top label lifts by assigning
top label to each coefficient. A zero polynomial with any label
lifts by placing that label at degree zero and $\tau$ elsewhere.
This proves surjectivity and compatibility with support restrictions.
Its exact fibre relation is equality of every amplitude coefficient
and equality of the total join; at a support stalk it is equality
of the join in $L_J$.

Diagonal entries in CLS1 retain each input label at its own degree,
while other entries contribute labels below existing ones. Therefore
the total join is unchanged. The amplitude calculation is the
ordinary polynomial exponential, giving
\begin{equation}\tag{CLS8}
 \eta_NF_{a,b}=
 \left[e^{-\operatorname{amp}(a)\partial_Z^2+
                   \operatorname{amp}(b)\partial_Z}\right]^\uparrow
 \eta_N,\qquad \eta_NP=\eta_NQ=\eta_N.
\end{equation}
On the common-label target the displayed linear operator keeps its
single label, and on support it is the identity. This is the exact
map under which the different projectors become equal.

\subsection{Synchronization on every divisor-indexed coefficient source}
Fix $d\ge2$ and the monomial basis of total degree at most $N$
in $s_1,\ldots,s_d$, with independent copies of $\mathcal M_k$.
The cyclic map $\beta$ replaces $Z^j$ by
$(s_1+\cdots+s_d)^j$, with the literal multinomial coefficients.
All nonzero coefficients are fully supported and all structural
zeros are $\tau$. Since each $D_iD_j\beta=\beta\partial_Z^2$,
the ordinary amplitude covariance $C$ acts by
\begin{equation}\tag{CLS9}
 \exp\left(\tfrac14\sum_{i,j}C_{ij}D_iD_j\right)\beta
 =\beta\exp\left(\tfrac14(\mathbf1^{\mathsf T}C\mathbf1)
                                           \partial_Z^2\right).
\end{equation}
For the original three matrices retain all entries and factors:
\begin{equation}\tag{CLS10}
 C_{\mathrm{ind}}=\frac td I,
 \quad C_{\mathrm{sync}}=\frac t{d^2}\mathbf1\mathbf1^{\mathsf T},
 \quad C_{\mathrm{rel}}=\frac td
       \left(I-\frac{\mathbf1\mathbf1^{\mathsf T}}d\right).
\end{equation}
Lift only the independent off-diagonal structural zeros as $\tau$,
and every entry of the other two matrices as supported, as in the
incoming source. Their off-diagonal sums are $e$, not $\tau$.
For each of these specified lifted matrices $\widetilde C$, define
$\mathcal U_C^\uparrow=
\exp(\tfrac14\sum_{i,j}\widetilde C_{ij}D_iD_j)$ on the stated
total-degree source. The finite exponential uses the same fully
supported reciprocal factorials and free identity as CLS1.
The exact identity on the coefficient sheaf is therefore
\begin{equation}\tag{CLS11}
 \mathcal U_{\mathrm{sync}}^\uparrow
 \mathcal U_{\mathrm{rel}}^\uparrow
 =\mathcal U_{\mathrm{ind}}^\uparrow P_\Sigma,
 \qquad
 P_\Sigma=\exp\left(\frac e4
                         (D_1+\cdots+D_d)^2\right).
\end{equation}
Indeed the sum of the lifted generators is the sparse independent
generator plus the displayed supported-zero generator: on the diagonal
adding $e$ does not change an already supported entry, and off the
diagonal it supplies the missing supported cancellation. The finite
commuting exponential proof of CLS2 gives the product identity.
The matrix $P_\Sigma$ has the identity on its diagonal, $e$ at
every coordinatewise lowering of positive even total degree, and
$\tau$ elsewhere. These are exactly all terms of powers of the
displayed derivative; each admissible path has a nonzero rational
coefficient and hence the specified label. Composing two such paths
again gives precisely that relation, proving $P_\Sigma^2=P_\Sigma$.

The cyclic formula retains the projector:
\begin{equation}\tag{CLS12}
 P_\Sigma\beta=\beta P,
 \qquad\mathcal U_{\mathrm{rel}}^\uparrow\beta=\beta P.
\end{equation}
For a source monomial $Z^j$, $\beta$ has top-supported coefficients
at every monomial of total degree $j$. For any lower same-parity
monomial choose a degree-$j$ monomial above it coordinatewise.
It receives a supported contribution under $P_\Sigma$; no other
degree or parity can receive one. Its amplitude remains zero in
that lower degree. This proves the first equality entry by entry.
For the second, the relative amplitude generator annihilates
$\beta$ by $\mathbf1^{\mathsf T}C_{\mathrm{rel}}\mathbf1=0$,
while its fully supported entries supply the same admissible
degree-lowering labels. This proves the second equality with its
original support, including $t=0^\bullet$.

Localization of CLS11 retains a visible difference. On the input
supported only at $s_1s_2$, $P_\Sigma$ contributes $1$ to the
constant label at every support stalk. Sparse independent heat
has no such constant contribution, since no $D_i^2$ can act on
that monomial. The synchronized and relative product does have it,
as CLS11 states. The amplitude of this extra term is zero before
localization, and the ordinary group reflection sends $P_\Sigma$
to the identity. Thus it is present on the actual support stratum
even though it is absent in the original scalar amplitude equation.

Finally the rational-prime maps $\alpha_r$ of ULC17 commute with
every matrix here. Each matrix acts on formal degrees by complex
scalars and labels; $\alpha_r$ is complex linear on coefficients
and fixes the labels. Hence every matrix entry commutes, proving
\begin{equation}\tag{CLS13}
 \alpha_r^{N+1}F_{a,b}=F_{a,b}\alpha_r^{N+1},\quad
 \alpha_rP_\Sigma=P_\Sigma\alpha_r,\quad
 \alpha_r\beta=\beta\alpha_r,
 \qquad k\longmapsto k+d(r).
\end{equation}
These identities hold between the actual divisor-indexed targets,
with the original prime and archimedean scaling in CDV22--24.
They do not assert that finite polynomial differentiation descends
through an unspecified spectral quotient or changes the zeros of
the original zeta function. The receiving maps, the localized
projectors and the full mixed-prime amplitude cohomology are now
specified on the same coefficient family.
